\begin{afterword}
\summaryhead

The post returns to the clever arguer of ``The Bottom Line,'' who says only true things about a box but chooses which things to say. Each statement is evidence, so it seems you must update on it; yet an arguer with enough signs to choose from could then make you believe anything.

The answer is that you should condition not on the reported facts alone but on the fact that a speaker following some rule reported them. A coin that may favour heads or tails is flipped ten times, and the speaker says that flips 4, 6 and 9 came up heads. Depending on the speaker's rule, the odds for a heads-biased coin are 8 to 1, 1 to 16, or something else. The Monty Hall problem likewise depends on how the host chooses a door.

Most legal processes rely on two opposed advocates; for a two-sided question that is almost as good as one curious inquirer, but not for questions with many sides. Filtering is no excuse to dismiss evidence you dislike: if you already know your side's case, a contrary fact still counts. Only a one-sided argument heard for the first time calls for caution.

\responsehead

The post is correct on its main point, and it answers the question that ``The Bottom Line'' left open. A true statement from a selective speaker is still evidence, but what it shows depends on how the speaker chose it. The coin and the Monty Hall cases show this cleanly, and their numbers are right. The closing limit, that ``a blue stamp on box B is still evidence,'' keeps the idea from becoming an excuse to ignore every argument one dislikes.

The rule the post gives assumes that the listener knows the speaker's rule of selection. In the coin and door cases the listener does, and knows how many flips or doors there are. In the case the post opens with, a hired arguer and a box, the listener knows the arguer's aim but not how many signs exist. The post does not work that case through. The author did so in a comment the next day, where the listener's probability rises as the arguer talks and drops when the arguer stops. Without that analysis the post's advice for the realistic case is ``beware.''

The post's limit on the idea depends on an unstated assumption. The post claims that if you are annoyed by a contrary argument, you are familiar with the case. Annoyance shows that you care about a question, not that you know its strongest arguments, and a listener hearing one side for the first time, the case the post treats separately, may be annoyed too.

The remark about courts is overstated. Two opposed advocates are the theory of common-law systems, not of ``most legal processes.'' In much of the world the court investigates the facts itself, which is closer to the post's curious inquirer than to two clever arguers.

The post also repeats two smaller faults noted in earlier posts. The post uses parallel worlds where a proportion over many boxes would do. And the footnote defines a bit as a measure of information, while the text uses the word for the logarithm of a likelihood ratio.

In short: a correct account of filtered evidence that answers the question ``The Bottom Line'' left open, though it does not work through the case it opens with, in which the listener does not know what the arguer left out.
\end{afterword}
